\documentclass[../main.tex]{subfiles}
\begin{document}
%==========================================TIÊU ĐỀ TẬP========================================
\begin{table}[h]
  \begin{adjustwidth}{-1cm}{}
    \begin{tabular}{>{\centering\arraybackslash}p{6cm}|>{\raggedright\arraybackslash}p{12.5cm}}
      \multirow{5}{6cm}
      % Cột bên trái: ÉPISODE và số tập
      {\\[-40pt]
        \flaregothic\fontsize{20pt}{20pt}\selectfont \centering
        \color{eptype}ÉPISODE\color{black}\\
        \burbank\fontsize{72pt}{72pt}\selectfont
        \color{epnum}47\color{black}\\[6pt]
        \cafeteria\fontsize{20pt}{20pt}\selectfont\color{epnum}(4-5)\color{black}
        \\[18pt]
        % \flaregothic\fontsize{14pt}{14pt}\selectfont
        % \color{epattr}BIENVENUE, \uline{\color{Banpire} Mel} !
        \flaregothic\fontsize{18pt}{18pt}\selectfont
        \color{epattr}ÉPISODE PERDU
        \color{black}
      }
       &
      % Dòng thứ nhất: tên tập tiếng Nhật
      {\fotskip\fontsize{18pt}{18pt}\selectfont \ruby{医}{い}\ruby{者}{しゃ}\ \ \ruby{vs.}{ヴァーサス}\ \ \ruby{DJ}{ディージェイ}
      } \\[6pt]
      % Dòng thứ hai: tên tập tiếng Anh
       & {\cafeteria\fontsize{18pt}{18pt}\selectfont Dr. Ayame}                                                  \\[4pt]
      % Dòng thứ ba: tên tập tiếng Việt
       & {\vnmsans\fontsize{18pt}{18pt}\selectfont Hội chứng \ruby{Sazae}{Sa-da-ê}-san\footnotemark}             \\[8pt]
      % Dòng thứ tư: Nhân vật xuất hiện
       & {\brian{\fontsize{14pt}{14pt}\selectfont Track used: Lost City - Celebrawl}
         }
      \\[8pt]
      % Dòng cuối cùng: Thời gian diễn ra của tập (thường sẽ là ngày phát hành)
       & {\continuum\fontsize{16pt}{16pt}\selectfont Ngày phát hành: 29/03/2020}                                 \\[18pt]
    \end{tabular}
  \end{adjustwidth}
\end{table}
\footnotetext{Còn được gọi là hội chứng ngày thứ Hai uể oải.}
%=========================================TRUYỆN CHÍNH========================================
\color{black}

\txtbx[2]{{\color{vol} \uline{Tóm tắt:}} Một buổi chiều Chủ Nhật yên ả bỗng chốc trở nên náo loạn bởi trò chơi khám bệnh tấu hài của Ayame và Pekora. Dựa vào một cuốn sách y khoa ``ảo ma'' mượn từ Fubuki, Ayame tự tin chẩn đoán Roboco mắc căn bệnh ``ảo tưởng DJ''. Sự thật ngã ngửa phía sau nhịp tim rock'n'roll của nàng Người máy lại chính là trợ lý A-chan đang kiệt sức nằm vất vưởng bên trong vì làm việc quá độ. Nhờ sự can thiệp kịp thời của bác sĩ Khan (bạn của Kuon), A-chan đã được cứu chữa và mọi người có được một bài học quý giá về sức khỏe.}

Một ngày Chủ Nhật cuối tháng Ba.

Trời xuân dịu dàng, gió nhẹ thoảng qua cửa sổ làm không khí trong căn phòng trở nên dễ chịu. Kuon đang ngồi đọc báo trên ghế sô-pha, thỉnh thoảng liếc mắt về phía Ayame và Pekora đang chơi trò ``khám bệnh'' đầy trẻ con.

``\hl{Bọn kia đang làm cái gì thế nhỉ? Nhìn giống trò con nít quá\dots}'' cậu buột miệng nói. Tiếng nói của cậu khiến Ayame, người đang đeo chiếc băng đô đồ chơi làm bác sĩ, liếc nhìn rồi phẩy tay: ``Anh cứ đọc báo của anh đi. Em đang bận chữa bệnh cứu người đây!''

Cậu Tổng nhướng mày nhìn bộ dạng của nàng Oni, nhưng rồi cũng gật đầu: ``À, vậy chắc anh không cần lo rồi\dots\ Chơi vui nha.''

Ayame ngay sau đó đổi giọng nghiêm nghị, nhìn về phía Pekora: ``Được rồi, há mồm ra nào, da-zo.''

\img{4-5a}

Nàng Thỏ làm giống như Ojou-san vừa nói. Sau khi nhìn bên trong cổ họng một cách đầy qua loa, cô nàng ``Bác sĩ'' kết luận: ``Hmm\dots\ Chị e là em đang bị mắc hội chứng S*zae-san rồi.''

Nàng Thỏ nhíu mày nhìn tiền bối cô, tỏ vẻ nghi ngờ: ``Thế là thế nào vậy, peko?''

Ayame đáp lại một cách đầy nghiêm túc như một bác sĩ thực thụ, nhưng không kém phần hài hước: ``Đừng có coi nhẹ căn bệnh đó nghen! Nếu bệnh tiến triển xấu đi, nó sẽ làm cho em đau đớn đến quằn quại đấy!''

Nghe chẩn đoán từ ``bác sĩ,'' Pekora nheo mắt, tỏ vẻ nghi ngờ với lời nói của nàng Quỷ sừng: ``Chị nghĩ Pekora này giống như một con ngốc lơ ngơ không biết gì chắc, peko?''

\img{4-5b}

Thấy cuộc tranh luận có vẻ thú vị, Kuon đặt điện thoại xuống bàn, rồi đi về phía hai người họ, tham gia vào cuộc trò chuyện: ``Anh không biết phải nói thế nào, nhưng mà em ấy nói đúng đấy.''

Nàng Thỏ quay ngoắt sang cậu, vẻ mặt không mấy tin tưởng với những lời thăm khám từ ``bác sĩ'' Oni kia: ``Cả anh cũng thế hả, Kuon-senpai!?''

Cậu gật đầu một cách bình thản: ``Từng bị rồi nên anh biết mà. Chỉ là theo cách khác hoàn toàn so với em nghĩ thôi.''

Lời giải thích của cậu khiến cho Pekora há hốc miệng ngạc nhiên. Cô không ngờ rằng một người nghiêm túc như cậu lại có thể nói ra những câu như vậy.

Trong khi cả hai đang tranh luận, Roboco xuất hiện, ghé đầu lại gần họ. Cô giơ tay chào ba người, hỏi: ``Có chuyện gì mà náo nhiệt thế\~{}?''

\img{4-5c}

Pekora nghe vậy, hớn hở đáp: ``Bọn em đang chơi trò khám bệnh, peko!'' Nàng Người máy thấy vậy, tỏ vẻ tò mò: ``Nghe thú vị nhỉ. Bác sĩ Ayame à, khám cho chị với được không\~{}?''

Nàng Quỷ sừng không hề do dự mà đồng ý ngay. Cô đưa tay lên cằm, làm điệu bộ như đang suy nghĩ: ``Hmm\dots\ Để xem nào\dots\ Chị nghĩ em bị thiếu dầu rồi đấy, Roboco-senpai!'' Cô nàng đọc kết luận ngay sau khi nhìn sơ qua cả cơ thể của nàng Rô-bốt.

Roboco phì cười, biết rõ là nàng Oni đang đùa: ``Thiếu dầu? Thế thì chị phải làm gì đây?''

Ayame khoanh tay lại, kê đơn cho Roboco y như một bác sĩ thực thụ: ``Đơn giản thôi. Chị về nhà bơm thêm dầu vào người nữa là khỏi ngay ấy mà, yo!''

Kuon nhìn nàng Quỷ sừng đọc vanh vách triệu chứng bệnh và cách chữa đầy tự tin liền bật cười nhẹ: ``Ayame-chan, làm gì ó chuyện khám ba giây là ra bệnh đâu. Với cả `hội chứng S*zae-san' chỉ là tên gọi chứ không phải bệnh thật đâu\dots''

Lập tức, nàng Quỷ sừng phồng má lên, tỏ vẻ phản đối: ``Ơ! Nhưng em thấy nó hợp lý mà! Với cả, nó thực sự là tên một loại bệnh đấy!''

Ngay sau đó, nàng Quỷ chuyển sang phấn khích, mắt sáng rực như vừa tìm được điều gì thú vị: ``Mọi người không biết đâu, nhưng em vừa đi tu luyện về và đã trở thành thần đồng trong lĩnh vực y khoa rồi đấy!''

Trong tay của nàng Oni là một cuốn sách bìa xanh lam, với bóng dáng của nàng Cáo trắng và một ô thoại ghi những dòng chữ: ``Những bài thuốc đơn giản cho Neko-conyan.''

\img{4-5d}

Pekora quay sang giải thích với Roboco, giọng nửa nghiêm túc, nửa chọc ghẹo với đôi tai thỏ rung rinh: ``Ít nhất là theo những gì chị ấy nói, peko.'' Kuon liếc qua cuốn sách rồi khẽ lắc đầu: ``Ờm\dots\ anh cũng không nghĩ em đang nghiên cứu đúng loại sách đâu\dots''

``Suy cho cùng ấy,'' Kuon thở dài, quay về phía nàng Roboco và nói: ``anh thấy Ayame-chan dường như chỉ đang chơi trò trẻ con thôi. Tin làm gì vào mấy cái ấy.''

Bất chợt, Roboco reo lên, ánh mắt đầy hào hứng khi hướng về phía Ayame và mở rộng vòng tay: ``Khám cho chị lại lần nữa xem nào\~{}! Lần này khám kỹ vào đấy.''

Cậu Tổng rướn mày, ngả người dựa vào ghế gần đó: ``\dots Em thực sự ngây ngô đến mức tin vào những gì em ấy nói sao, RoboPON-chan?''

Đến cả nàng ``Bác sĩ'' tự xưng cũng phải nhắc khéo: ``Nhưng chị là rô-bốt mà. Vừa nãy những gì em nói là dựa theo bản tính hậu đậu của chị mà em đoán đó, yo.''

Mặc cho hai người châm chọc cô đến mức Roboco phồng má lên, cô nàng vẫn tỏ vẻ kiên quyết không chịu rời đi. Nàng Quỷ sừng cũng chỉ đành miễn cưỡng chấp nhận lời đề nghị khám lại của cô ấy.

Cô cầm chiếc ống nghe đồ chơi, đặt nhẹ lên ngực trái của nàng Người máy. Ban đầu, mọi thứ có vẻ bình thường, nhưng rồi\dots\ dường như có một bản nhạc rock bùng cháy đang phát ra.

\img{4-5e}

Ayame giật mình, lùi ra phía sau, mặt tái mét lại, mắt trợn ra tỏ vẻ sợ hãi, không thốt nổi một từ nào. Roboco giật thót, tỏ vẻ lo lắng: ``C-Chị bị mắc bệnh gì sao?''

Pekora không thể giữ được sự tò mò, lập tức nhích lại gần nàng Quỷ, mắt cô sáng như sao một cách kỳ lạ: ``Sao, sao? Chị nghe được gì thế, peko!?''

Không chờ nàng Oni trả lời, cô nhanh chóng giật lấy chiếc ống nghe từ tay nàng Quỷ và áp vào ngực của ``bệnh nhân'' một lần nữa. Lần này, một kiểu nhạc khác lại phát ra, nghe giống như một kiểu điệu nhạc đến từ thế giới Ả rập được chơi trên đàn oóc-gan.

Nàng Thỏ nhíu mày, rồi thốt lên với vẻ tỉnh bơ: ``Nghe giống như mấy cái bài thường nghe trong siêu thị gần nhà em vậy.''

\img{4-5f}

Nàng Oni chống tay vào hông, hỏi nàng Thỏ với sự khó hiểu: ``Bộ siêu thị gần nhà em chơi nhạc rock'n'roll sao?''

``Không có đâu, peko!'' nàng Thỏ đáp lại đầy quả quyết.

Đến lượt cậu con trai cũng hiếu kỳ với những gì họ nghe thấy: ``Mấy đứa nói thế nào chứ\dots\ Đưa anh nghe thử xem nào.''

Kuon lấy chiếc ống nghe từ tay nàng Thỏ và áp vào ngực của Roboco. Ban đầu, khuôn mặt cậu vẫn bình thường, nhưng chỉ vài giây sau, đôi chân của Kuon bất ngờ bắt đầu nhảy nhót theo nhịp điệu, miệng còn lẩm nhẩm kêu ``beat, beat,'' thậm chí còn vung tay lên trời đầy nhiệt huyết.

Nàng Quỷ sừng cau mày, tỏ vẻ bất ngờ với những gì cậu nghe được: ``Anh nghe được gì mà trông sung thế?''

Kuon vừa cười, vừa hơi lắc lư theo nhịp của giai điệu, trả lời một cách đầy hào hứng ``Anh lại nghe giống như mấy cái nhạc ở hội chợ lúc đi chơi với mẹ anh vậy. Nghe đã lắm!''

Trong khi ba người đang sôi nổi bàn luận về những âm thanh kỳ quặc bên trong ngực của nàng Người máy, bản thân cô nàng đứng cạnh cảm giác mình bị bỏ rơi, cuối cùng không chịu được nữa, cô bực bội lên tiếng, tỏ vẻ mất kiên nhẫn: ``Cho chị nghe thử với\~{}!''

Roboco nhanh chóng cầm lấy ống nghe từ tay Kuon, áp vào ngực mình. Nhưng không như kỳ vọng, thứ phát ra từ bên trong cơ thể cô không phải âm nhạc, mà là một giọng nói kỳ lạ, trầm thấp như thể đến từ người bí ẩn: ``\dots chưa? Mọi người đã cảm thấy sức mạnh chưa?''

Nàng Người máy rút ống nghe ra ngay lập tức, giọng run rẩy, ánh mắt lộ rõ vẻ hoang mang: ``Con người xấu trong người chị đang gọi chị nè\~{}\dots''

\img{4-5g}

Nhìn ánh mắt lo âu của cô, Kuon tỏ vẻ ngạc nhiên: ``Ủa? Thế là thế nào? Sao mỗi người nghe thành một kiểu thế?''

Nàng ``Bác sĩ'' bối rối gãi đầu, rồi giở giọng tỏ vẻ nghiêm trọng hơn hẳn: ``Triệu chứng này là gì không biết nữa!? Để em tra thử.''

Nàng Oni lập tức lật cuốn sách y khoa đầy tính "nghiên cứu" của mình. Trong khi đó, Pekora và Kuon tò mò đứng bên cạnh nhìn vào những dòng chữ được ghi chép một cách ngăn nắp nhưng lạ kỳ.

Sau một hồi lật giở, cả ba người dừng lại ở một trang có ghi tên bệnh là: ``Ảo tưởng NVC trong trò chơi thám hiểm.'' Nói tóm gọn lại: ``Hội chứng tuổi mới lớn'' hay ``chuunibyou.''

Một tong những dòng mô tả của bệnh làm cả ba phải cau mày: ``Có những hành động làm cho người bệnh hét vào gối vào ban đêm khi họ già đi.'

Cả ba người họ nhìn nhau, rồi cùng lúc quay đầu nhìn Roboco, nở một nụ cười mang ý châm chọc, dường như muốn nói: ``Chắc chắn là Roboco-senpai/-chan mắc căn bệnh này rồi!''

\begin{figure}[H]
  \centering
  \begin{subfigure}[t]{0.45\textwidth}
    \centering
    \includegraphics[width=8cm]{../../../../Image/Book4/4-5h1.png}
    \caption{Một đoạn giải thích triệu chứng của hội chứng chuunibyou.}
  \end{subfigure}
  \quad
  \begin{subfigure}[t]{0.45\textwidth}
    \centering
    \includegraphics[width=8cm]{../../../../Image/Book4/4-5h2.png}
    \caption{Ayame, Pekora và Kuon (giả sử cậu ở đằng sau hai cô gái) đọc triệu chứng, vẻ mặt tỏ rõ sự ngỡ ngàng.}
  \end{subfigure}
  \quad
  \begin{subfigure}[t]{0.45\textwidth}
    \centering
    \includegraphics[width=8cm]{../../../../Image/Book4/4-5h3.png}
    \caption{Họ nhìn về phía Roboco, tin rằng cô nàng đang bị hội chứng tuổi mới lớn kia.}
  \end{subfigure}
\end{figure}

Mặt của nàng Người máy bắt đầu tái dần, giọng lo lắng xen lẫn cầu cứu: ``C-Cái ánh mắt kia là sao vậy!?''

Kuon bật cười, nhìn tên bệnh trong cuốn sách với vẻ thích thú: ``Có gì đâu. Bọn anh đang nghi là em đang mắc bệnh\dots''

Chưa kịp để cậu Tổng nói hết câu, một cảnh tượng kỳ quái xảy ra. Từ ngực của Roboco, hai bàn tay thò ra như thể ai đó đang cố trồi lên từ bên trong. Pekora giật mình, hoảng hốt hét lớn: ``\textbf{Peko!? Cái gì thế này!?}''

Cả nhóm nhìn chằm chằm vào "sinh vật" đang xuất hiện, và không ai khác, người vừa trồi lên lại chính là A-chan. Cô nàng trợ lý, với vẻ mặt không hề hốt hoảng, thay vào đó đáp lại một cách đầy phấn khích (dù khuôn mặt mệt mỏi của cô không thể hiện ra): ``Tôi hỏi là mọi người có cảm thấy sức mạnh hay không á!''

Kuon đứng hình trong vài giây, lắp bắp khi nhìn thấy đồng nghiệp của mình đang ở bên trong cơ thể của Roboco: ``\dots A-san!?''

Pekora trố mắt nhìn, thốt lên với vẻ ngạc nhiên: ``A-chan!? Sao chị lại ở đó, peko!?''

Roboco thì ôm đầu hoảng sợ, lắc qua lắc lại như không tin vào những gì đang xảy ra. Trong khi đó, Ayame vẫn bình tĩnh hơn cả, tiếp tục lật cuốn sách với vẻ tìm tòi.

Sau một lúc tra cứu, cô nàng ngẩng đầu, rồi thốt lên như vừa tìm ra đáp án: ``Không lẽ là\dots? Đây là căn bệnh `ảo tưởng DJ\footnote{Nguyên văn: DJ病 (diijeibyou), tức là căn bệnh DJ. Để vần với 中二病 (chuunibyou), người viết sẽ để là ``ảo tưởng DJ''.}' mà!''

Nàng Thỏ và nàng Rô-bốt đồng thanh, kinh ngạc hét lên: ``Bệnh `thoả hiệp với một con quỷ DJ' ư!?'' Trong khi đó, Kuon nhíu mày, nhìn nàng Quỷ sừng như thể không tin vào tai mình: ``Sao lại có cả cái bệnh nghe lạ đời thế\dots?''

Nàng ``Bác sĩ'' giơ cuốn sách lên, chỉ vào một đoạn cụ thể, tỏ vẻ đầy chắc nịch: ``Thì trong này nó ghi như thế mà! Bệnh này làm người bệnh bị cuốn vào những bản DJ đến mức họ bắt đầu mơ thấy mình là một nghệ sĩ DJ, sau đó bị `quỷ DJ' (tức là lý trí) điều khiển hoàn toàn!''

Nghe triệu chứng bệnh mà nàng Oni đọc lên, Roboco càng tỏ vẻ hoảng loạn hơn: ``Chị\dots\ chị bị cái bệnh đó hả!?'' Thấy vậy, cậu Tổng lên tiếng trấn an: ``Anh không chắc là bệnh thật hay chỉ là do cái cuốn sách của em ấy viết như vậy nữa\dots''

A-chan, người vẫn đang lủng lẳng trước ngực của Roboco, tỉnh bơ nói: ``Bệnh quái gì chứ, Nakiri-san? Đây là sức mạnh đặc biệt chứ! Tôi chỉ muốn kiểm tra mọi người có đủ tố chất để tiếp nhận không thôi.''

\img{4-5i}

Nàng Oni hạ cuốn sách xuống, bĩu môi: ``Thế chẳng phải giống mấy con quỷ trong truyện sao? Dụ người ta xong mới lộ bản chất thật!''

Đáp lại, nàng Đeo kính chỉ nháy mắt cười, đáp lại đầy ẩn ý: ``Thì đôi khi cũng phải có chút kịch tính cho cuộc sống thêm thú vị mà, đúng không?''

Cả nhóm sững người trước lời giải thích vừa đúng vừa sai của cô nàng trợ lý. Pekora lùi lại, kéo tay Kuon, thì thầm: ``\hl{Chị ấy\dots{} bình thường có quái vậy đâu, peko?}''

Kuon lắc đầu: ``Có khi chị ấy bị ảo tưởng thật rồi.'' Ngay khi cậu dứt lời, khuôn mặt A-chan lập tức sầu não và buồn tủi. Cô than thở: ``Mọi người thật sự nghĩ tôi bị ảo tưởng ư? Nghe không hợp lý tẹo nào!''

Trong khi đó, Roboco vẫn ngúng nguẩy, cố gắng giữ thăng bằng trên thân mình để không ngã. Cô ``Bệnh nhân bất đắc dĩ'' đáp lại, vẻ mặt đầy bối rối:
``Thì chị gần đây hay làm các chương trình DJ thật mà.''

Pekora nghiêng đầu, thắc mắc với vẻ ngây thơ: ``Vậy là bọn em nghe thấy nhạc vì chị đang DJ (chà đĩa nhạc) à, peko?''

Nàng ``Tăng ca'' khoanh tay, phồng má giận dỗi, than vãn với cái tên bệnh trong cuốn sách ngớ ngẩn mà nàng Oni đang cầm kia: ``Kiếm cho chị cái tên nào dễ nghe hơn coi\dots''

Kuon nhướng mày, nhìn cô nàng đang đung đưa theo nàng Người máy đó, tò mò hỏi: ``Mà sao chị lại chui tọt vào ngực của Roboco-chan để chơi DJ vậy?''

A-chan đáp lại bằng một tông giọng thản nhiên, như thể đó là điều bình thường: ``Vì nó yên tĩnh hơn đấy em.''

Nàng Quỷ Ayame, vẫn chưa thỏa mãn, tiếp tục tra hỏi người vừa ló ra từ thân của Roboco: ``Bình thường DJ-chan hay làm gì vậy?''

Cô nàng Đeo kính đáp lại một cách ráo hoảnh: ``Làm công việc văn phòng thôi.'' Cả Kuon và Pekora cùng đồng thanh, hét lên đầy ngao ngán: ``\textbf{Cái đó thì ai chả biết(, peko)!}''

\ct

Khi Roboco đã bình tĩnh trở lại sau một hồi hoảng loạn, cô quay sang hỏi ``Bác sĩ'' Ayame với vẻ đầy hy vọng: ``Có cách nào để chữa căn bệnh này không?''

Nàng Quỷ sừng lật cuốn sách một cách nghiêm túc, sau đó ngẩng đầu lên và nói với cả nhóm bằng giọng đầy kịch tính: ``Có vẻ như cô ấy sẽ quằn quại đến nổ tung trong đau đớn nếu ta đọc câu thần chú này.''

\img{4-5j}

Nhìn những dòng ghi cách chữa căn bệnh trên cuốn sách, cậu Tổng nghiêng đầu, tỏ vẻ đầy hoài nghi: ``Khoan\dots\ Ta chữa một căn bệnh\dots\ bằng cách mắc bệnh cho chị ấy sao?''

Mặc cho những câu hỏi vẫn còn vương vấn trên đầu mọi người vì cách chữa trị kỳ lạ này, họ đều phải công nhận rằng đây là cách chữa duy nhất cho căn bệnh có tên vừa dài ngoằng vừa oái oăm này.

Pekora hào hứng đề nghị: ``Vậy tất cả bọn mình đều cùng đọc nó ha. Cả anh nữa, Kuon-senpai!''

Kuon nghe vậy, thở dài, gãi đầu bất lực: ``Xin lỗi chị nhé, A-san. Bọn em thực sự không muốn làm việc này đâu\dots'

Họ lấy hơi chuẩn bị, đồng loạt gào lên: ``\textbf{Ngày mai là ngày thứ Hai rồi đó chị!}''

Ngay lập tức, ``cái u'' trên ngực Roboco (vốn là A-chan) kêu lên đầy đau đớn: ``Không! Không phải là thứ Hai chứ!''

Cô nàng tóc tím tạo dáng như một ngôi sao Rock'n'Roll, trước khi phát nổ (đúng nghĩa đen) một cách hài hước. Áp lực vụ ``nổ'' khiến A-chan tuột hẳn ra khỏi ngực của Roboco, lăn lông lốc xuống sàn.

\gif{4-5k}{1}{78}

Roboco ôm ngực, chưa tin vào mắt mình: ``Ơ\dots\ thế là\dots\ hết bệnh rồi đúng không?'' Trong khi đó, Kuon, Pekora và Ayame chỉ biết đứng nhìn A-chan đang lồm cồm bò dậy, vừa ho sặc vừa làu bàu: ``Mấy đứa\dots\ đùa quá trớn rồi đó\dots!''

Và rồi, có vẻ như cô nàng đã mất hết năng lượng sau câu nói vừa thô vừa thật này, cuối cùng nàng Đeo kính lăn đùng ra ngất xỉu.

%==========================================TRUYỆN PHỤ=========================================
\omk

Sau khi A-chan hoàn toàn rời khỏi người Roboco và ngất xỉu trên sàn sau cú hét mang tên ``hội chứng thứ Hai'' của cả nhóm, Ayame và Pekora vội lao đến bên nữ quản lý để tìm cách cấp cứu.

\img{4-5l}

Nàng Quỷ sừng lo lắng, vừa lay người bất tỉnh vừa thốt lên: ``Chị ấy ngất xỉu rồi! Ta phải làm gì đây!?''

Nàng Thỏ lắc đầu bối rối, giọng hơi run run: ``Pekora không có biết đâu, peko! Hay là thử hồi sức nhân tạo cho chị ấy?''

Trái ngược với sự hoảng hốt của hai cô nàng, Kuon đứng ở gần đó chỉ lắc đầu ngán ngẩm, bình tĩnh phán một câu: ``Chị ấy làm việc vất vả quá nên mới như vậy đấy.''

Roboco, sau khi cảm thấy ``căn bệnh'' của mình được chữa khỏi, ôm lấy ngực, thở phào nhẹ nhõm. Thấy tình hình có vẻ ổn với cô nàng, cậu con trai nhanh chóng lên tiếng chỉ đạo: ``\textbf{PON ơi!}''

Nghe thấy mình bị trêu chọc, cô nàng Rô-bốt liền phồng má, phản ứng gắt như mọi khi: ``Em bảo rồi, em không có PON đâu mà!''

Cậu Tổng nhún vai đầy thản nhiên, rồi tiếp tục chỉ đạo: ``Sao cũng được. Đưa chị ấy xuống phòng y tế đi. Bảo bác sĩ truyền nước muối vào rồi cho chị ấy nghỉ ngơi.''

Roboco liền chạy tới nàng Đeo kính, nghiêm túc đáp: ``Em biết rồi!'', rồi vác A-chan đi thẳng xuống phòng y tế. Ayame và Pekora nhìn theo với ánh mắt ngạc nhiên, rồi quay sang cậu con trai.

``Sao anh biết được vậy?'' nàng Quỷ sừng nghiêng đầu thắc mắc làm sao mà cậu có thể đoán được như vậy.

Cậu con trai nhìn nàng Người máy cõng cô nàng ngất xỉu kia, khoanh tay đáp lại: ``Có gì khó đâu. Chị ấy bị bệnh suy nhược cơ thể do làm việc quá sức.''

Dường như vỡ lẽ được tình huống đang xảy ra, nàng Thỏ gật gù: ``À\dots\ Em cũng từng nghe qua rồi, peko.'' Nhưng rồi, cô nàng tiếp tục thắc mắc: ``Nhưng mà, sao biết mệt thế mà chị ấy vẫn chơi DJ được nhỉ?''

Kuon bất lực, giơ hai tay lên trời: ``Bó tay.''

\ct

Ngay sau đó, Kuon liên lạc với Hijikara Khan để thăm khám sức khỏe của A-chan. Hiện cậu vừa là bác sĩ tại một bệnh viện tư vừa kiêm lập trình phần mềm máy tính. Ngay khi nhận được lời mời, Khan đã nhanh chóng có mặt và tiến hành khám chữa cho cô nàng Đeo kính kia - người cũng vừ mới tỉnh lại sau khi truyền nước.

Sau khi kiểm tra một lượt, cậu Bác sĩ (hàng thật) gật đầu kết luận: ``Cô ấy cần nghỉ ngơi nhiều hơn và bổ sung thêm dưỡng chất. Đây là triệu chứng điển hình của việc làm việc quá sức.''

Nghe đến hai chữ ``nghỉ ngơi'', A-chan dù đang nằm trên giường vẫn phản đối kịch liệt: ``Không được! Tôi không thể nghỉ ngơi được! Ai sẽ lo công việc nếu tôi không có ở đó?''

Kuon liền lên tiếng trấn an: ``Chị nghỉ đi. Phần việc của chị, để em lo.'' Cô nàng ``Tăng ca'' nhìn chàm chằm vào cậu, như muốn xác nhận lại: ``Cậu chắc chứ? Tôi vẫn còn nhiều việc phải làm đấy!''

Đáp lại cô là một cái gật đầu tự tin, đôi mắt ánh lên một sự nghiêm túc: ``Chắc như đinh đóng cột. Đằng nào mấy ngày nay em cũng đang rảnh.''

Cuối cùng, sau một hồi chần chừ, cô miễn cưỡng gật đầu nghe theo chỉ định của bác sĩ.

Ngay sau khi thăm khám và kê đơn cho cô, Khan quay sang nhóm ba người Pekora, Roboco, và Ayame với vẻ mặt nghiêm nghị: ``Ba người các cô cũng nên học qua vài kiến thức cơ bản về các loại bệnh, dấu hiệu nhận biết và cách phòng tránh. Sống trong thế giới loài người thì đây là những điều cần thiết.''

Nàng Quỷ sừng nghiêng đầu tò mò: ``Kiểu như bệnh `ảo tưởng DJ' hả?'' Cậu Bác sĩ khẽ nhếch môi cười, chỉnh lại cặp kính: ``Không, đó là bệnh lý\dots\ tận đẩu tận đâu ấy. Tôi nói về những thứ thực tế hơn, như cảm cúm, suy nhược, thiếu ngủ chẳng hạn.''

Nàng Thỏ gật gù như đã hiểu, nhưng mắt vẫn ánh lên vẻ tò mò: ``Nghe có vẻ phức tạp quá, peko! Anh có thể nói rõ hơn được không?''

Roboco chớp mắt, hỏi một cách hồn nhiên: ``Nhưng nếu em không bị cảm, thì có cần học không?''

Khan thở dài, vỗ trán: ``Tôi biết cô là rô-bốt nên không thể cảm thấy bị bệnh được, nhưng rồi sẽ có lúc cần đến, tin tôi đi.''

\ct

Sau khi A-chan được đưa về phòng nghỉ ngơi, Kuon bước ra khỏi phòng y tế với dáng vẻ thoải mái hơn. Cậu quay lại nhìn người bạn cũ của mình, rồi cảm ơn: ``\vie{Cảm ơn cậu nhé, Khan. Có lẽ không ai có thể làm tốt hơn cậu trong khoản chữa bệnh này.}''

Khan vỗ vai vào cậu Tổng, nở một nụ cười nhẹ nhàng: ``\vie{Có gì đâu, bạn cũ gặp lại mà. Nhưng mà này, Kuon\dots\ lần sau hãy giữ gìn sức khỏe cho chị ấy tốt hơn đi. Làm tổng quản lý mà lại để đồng nghiệp làm việc kiệt sức như vậy thì không ổn đâu.}''

Kuon khẽ gật đầu, cười trừ nhìn về phía A-chan đã say giấc nồng: ``\vie{Ừ, việc này để tôi ý kiến lên ông sếp.}''

Trong lúc hai người nói chuyện, nhóm Pekora, Roboco, và Ayame cũng vừa đi tới. Ayame vẫn không ngừng lật giở cuốn sách về các loại bệnh: ``Bác sĩ Khan này, thế mấy cái bệnh như `ảo tưởng DJ' có cách chữa không?''

Cậu Bác sĩ quay mặt lại nhìn cuốn sách mà cô nàng Oni đang cầm trên tay. Cậu đáp lại với một tông giọng vừa đùa vừa thật: ``Cách chữa đơn giản nhất là đừng bao giờ nghĩ đến chúng nữa. Cô kiếm cái cuốn sách này ở đâu vậy?''

Nghe câu trả lời từ Khan, nàng Thỏ bật cười, lém lỉnh nói: ``Vậy chắc Ayame-senpai không chữa được đâu, peko!'' Ayame lườm cô, phồng má lên giận dỗi: ``Này! Đừng có trêu chị chứ! Mất công lắm chị mới mượn được từ Fubuki-senpai đó!''

Kuon nhìn liếc qua bìa của cuốn sách đó, lúc này mới nhận ra hình bóng của nàng Cáo trắng. Cậu thở dài, cười chát: ``Hèn gì. Mượn ai không mượn lại đi mượn đúng cái tên pha lắm trò nhất trong công ty\dots''

Roboco khẽ cười với câu nói đùa của cậu, nhưng rồi ngập ngừng hỏi: ``Mà\dots\ A-chan sẽ ổn chứ? Em thấy chị ấy với cả senpai làm việc cực kỳ nhiều mà.''

Cậu Tổng khẽ thở dài, xoa trán: ``Anh thì mấy đứa không phải lo. Còn chị ấy thì sẽ ổn thôi. Chị ấy chỉ cần nghỉ ngơi một thời gian là khỏe lại. Nhưng anh mong rằng lần này sẽ là bài học cho chị ấy\dots\ và cả mọi người nữa.''

Cậu liếc nhìn về cơ thể máy móc của cô, buông ra một lời châm chọc: ``Cơ thể của con người bọn anh không phải là rô-bốt như của em đâu, RoboPON-chan.''

``Lại gọi em là PON nữa! Mồ!'' cô đỏ mặt, phồng má phụng phịu tỏ vẻ phản đối.

Tiếng cười vang lên giữa nhóm bạn. Mặt trời bên ngoài khẽ lặn, ánh chiều tà nhuộm màu vàng cam lên khung cảnh xung quanh.

Khan nhìn chiếc đồng hồ đeo tay, rồi quay sang chào tạm biệt: ``\vie{Thôi, tôi về trước đây. Hôm nay bệnh nhân của ông làm tôi cũng mệt phát ốm!}''

Kuon cười nhạt, vẫy tay: ``\vie{Ừ, cảm ơn nhé. Hôm nào ông rảnh thì để tôi khao ông một bữa.}'' Cậu Bác sĩ quay đầu lại, giơ tay làm dấu ``OK'': ``\vie{Nhớ đấy, bạn hiền!}''

Khi bóng của Khan khuất dần, nhóm bốn người còn lại đứng nhìn nhau. Ayame nheo mắt, hỏi: ``Vậy, bây giờ làm gì tiếp đây, Kuon-senpai?''

Pekora bật cười tinh nghịch: ``Chắc mọi người cũng đói rồi, hay là ta ăn tối đi, peko!''

Cậu Tổng nhìn đồng hồ đeo tay, rồi thở dài: ``Ừ, anh cũng cần làm nốt vài thứ cho A-san trước khi kết thúc ngày hôm nay.''

Nàng Người máy cười nhẹ nhàng: ``Vậy thì đi ăn thôi. Em đang đói rồi.''

Dưới bầu trời chiều nhuộm sắc vàng ấm áp, cả nhóm chầm chậm bước đi, tiếng nói cười rộn ràng vang khắp hành lang. Một ngày dài đầy hỗn loạn cuối cùng cũng khép lại, để lại dư vị nhẹ nhàng và ấm áp của tình bạn, sự quan tâm lẫn nhau giữa những con người\dots\ và rô-bốt.

\end{document}